\section{A revised fact changes a conclusion without rewriting its history}
\label{sec:sector-evidence-revision}

A correction can change the amount that a model now explains while leaving
both its earlier explanation and the company's committed journals intact.
The bank and insurance examples later in this proposal make the distinction
concrete. The bank needs the credit-risk baseline of its irrevocable
commitment even after a loan is drawn. The insurer needs to know whether
its reinsurance treaty preceded the recognition of an onerous group.
Neither fact can be replaced by the time a document happens to arrive.

Consider a constructed bank record that supplies no original commitment
baseline. The analyst can retain its cash flows and conditional loss
scenarios, but cannot establish the required original-to-current credit
comparison. A later record identifies the original commitment risk and
supports the separately supplied finding of significantly increased credit
risk. With the other assumptions in Section~\ref{sec:bank-commitment-draw},
the funded allowance is 4 and equity is 16. The additional record resolves
the missing history; it does not reset the baseline to the draw date.
The relevant AASB~9 provisions retain the irrevocable commitment as the
initial comparison and treat the resulting asset as its continuation
\citep[paragraph 5.5.6 and B5.5.47]{aasb9dated}.
The probability tree remains an authored assumption, and the credit finding
is not inferred merely because a probability doubled.

Suppose, for this constructed record, exact bytes were publicly available
on 2 January, ingested into the system on 5 January, and accepted for the
required reading on 6 January. An investor-information reconstruction with
a 3 January cutoff may use those bytes if their public availability is
verified, even though this particular system obtained them later. A
system-as-of query for 3 January cannot use them. On 5 January a raw
system query may retrieve the document, but a query for an accepted reading
still lacks the 6 January decision. Later extraction is compatible with
the earlier investor cutoff only if its interpretation introduces no later
information. These dates are authored specification fixtures; they
establish no actual historical availability.

\subsection*{A treaty correction separates knowledge from measurement}

In Section~\ref{sec:insurance-connected-groups}, treaty R entered before
group B's recognition produces an initial reinsurance asset of 85,
initial equity of 85 and year-one equity of 83. If R instead enters
immediately afterward, with all other stated facts unchanged, the initial
asset is 70, initial equity is 70 and year-one equity is 74.75.
The loss-recovery recognition condition causes the difference
\citep[paragraph 66A and B119C]{aasb17}. Both histories eventually
reach cash of 88.3 on the same supplied claim and recovery path.
Agreement on that terminal amount therefore cannot establish correct
recognition timing.

Now give the two order assertions different information dates. A constructed
January disclosure says the treaty preceded recognition; a February correction
says it followed recognition and explicitly replaces that order assertion.
A latest-available query cut off in January retains the January assertion.
The February record is known to be too late and cannot suppress it.
A February query may select the correction and produce the later-entry
measurement. The correction concerns the order assertion for B and R;
it does not replace every other assertion in the company's report.

A competing record with unknown public availability has a different effect
from a record known to be published after the cutoff. The known late record
is excluded. The unknown record leaves an unresolved dependency if it might
affect the requested conclusion, even while an older eligible assertion is
displayed. Similarly, an asserted publication date without evidence linking
the retained bytes to that date does not establish historical admission.
If only the common calendar date of treaty entry and group recognition is
known, their order remains unresolved; the model cannot choose 85 merely
because that treatment is more favorable.

An extraction correction is different again. A source may consistently say
that R followed recognition while its first accepted extraction mistakenly
records the opposite. Correcting the reading creates a new reading and
review decision linked to the same source bytes. The old result remains
retrievable as a historical error; reproducing it does not make it correct.
An issuer's correction instead changes the source assertion and usually
its bytes. The system records which change occurred before deciding which
derived results require reassessment.

\subsection*{Invalidate the selection, preserve the committed event}

A result depends on the candidates considered by its selector as well as
the candidate ultimately chosen. Adding a possible competing treaty record
changes that candidate population. Reusing a cached selection solely
because its chosen assertion identifier is unchanged would miss the new
obligation. The proposed dependency graph therefore includes a version for
candidate-set membership. The change propagates from membership through
selection and interpretation to measurement and the resulting statement.
A fresh selection may exclude a known late competitor and calculate the
same amount, but it records that decision under a new view.

Newly created descendants must check their complete ancestry too. Otherwise
a report created after an invalidation scan could inherit a stale measurement
without appearing on the scan's earlier list. In contrast, a bank result whose
complete dependency set contains no changed insurance inputs can remain
reusable. Neither a shared company label nor an unchanged displayed value
establishes independence. Missing ancestors leave reuse unestablished; a
cycle or missing required reference rejects the affected dependency chain.

The retained explanation of a result uses its original evidence view.
It can show a current warning that a dependency has changed, but it cannot
replace the old source with today's correction and still purport to explain
the old calculation. Requesting a different view requires a new calculation.
This preserves the distinction between the company's historical record,
what the system then concluded, and what a later reconstruction supports.

The same rule protects committed accounting events. The connected insurance
example commits financing, service, incurrence and settlement as four
separate events. After settlement, retrying the exact original financing
request returns its original receipt and reports no additional commit;
today's evidence revision does not turn that retry into a new assessment.
Changing the request while retaining the business-event identity is a
conflict. The request's original bindings, historical admission and required
bytes must still validate before receipt lookup.

A new conclusion that the treaty came later does not itself authorize an
entry reducing the committed reinsurance asset. A correction proposal must
identify its accounting treatment, effective period, affected accounts and
opening-balance consequences under the applicable authority. Without that
support, the proposed assessment returns \code{TREATMENT\_UNSUPPORTED},
retains the revised conditional calculation and creates no journal.
Versioning supplies the record of a correction; it cannot supply the
accounting rule that makes the correction appropriate.
